\bigskip\Needspace{10\baselineskip}
\pdfbookmark[1]{Ceres Ferdinandea. November 1802}{ceres180211}
\begin{center}
CERES FERDINANDEA.\par
\medskip\hrule\medskip
Monatliche Correspondenz. Vol. VI. pp.\ 492—498. November 1802.\par
\medskip\hrule
\end{center}\medskip
—\quad —\quad —

Dr.\ \textsc{Gauss} has meanwhile continued working on the orbit of this planet. Taking into account the perturbations as he calculated them (October issue, p.\ 387) [p.\ 224 of this volume] from his elements VII, together with the latitude equations, for which his calculation and \textsc{Oriani}’s (June issue, p.\ 586. July issue, p.\ 68) [1802. July pp.\ 67—70. —\quad —\quad —\quad \textsc{Oriani}, a skilled calculator so well versed in intricate perturbation calculations, took the trouble to repeat the entire perturbation calculation for \textit{Ceres} with a different assumption about the mean distance and thereby check his original results again. Here is the outcome of this entire calculation, which this great astronomer kindly shared with us.
\[
\begin{aligned}
\text{Let }D&=\text{mean longitude of }\CeresSymbol-\text{mean longitude of }\JupiterSymbol;\\
A'&=\text{mean anomaly of }\JupiterSymbol;\\
A&=\text{mean anomaly of }\CeresSymbol.\\
H'&=\text{mean longitude of }\JupiterSymbol-\text{mean longitude of }\AscendingNode\JupiterSymbol.\\
H&=\text{mean longitude of }\CeresSymbol-\text{mean longitude of }\AscendingNode\CeresSymbol.
\end{aligned}
\]
It should also be noted that the perturbation equations for any other eccentricity $e$ of the orbit are obtained by multiplying the terms containing $A$ by $\dfrac{e}{0.081406}$, the terms containing $2A$ by $\left(\dfrac{e}{0.081406}\right)^2$, the terms containing $3A$. by $\left(\dfrac{e}{0.081406}\right)^3$.

\Needspace{24\baselineskip}
The resulting equations for the perturbations of \textit{Ceres} by Jupiter, with their changes, are given below.
\SourceLeaf{233}%
The following table gives the perturbations of the heliocentric latitude.
\begin{center}
\begingroup\small
\begin{tabular}{l|r}
\shortstack{According to \textsc{Gauss}’s\\orbital elements VII}&\shortstack{If the mean annual motion\\of Ceres is increased by\\20 min.}\\\hline
$-13''19\sin(H-D)$&$-12''97$\\
$+1.62\sin H'$&$+1.59$\\[3pt]
$+16.21\sin(2D-H)$&$+16.03$\\
$-1.99\sin(D-H')$&$-1.96$\\[3pt]
$+32.38\sin(3D-H)$&$+32.78$\\
$-3.97\sin(2D-H')$&$-4.01$\\[3pt]
$-5.50\sin(4D-H)$&$-5.31$\\
$+0.67\sin(3D-H')$&$+0.65$\\[3pt]
$+6.17\sin(D+H)$&$+6.06$\\
$-0.76\sin(2D+H')$&$-0.74$\\[3pt]
$+1.52\sin(2D+H)$&$+1.49$\\
$-0.19\sin(5D+H')$&$-0.18$\\[3pt]
$+16.15\sin(4D-A'-H-2^0\ 33')$&$+13.83$\\
$-13.86\sin(5D-A-H-2^0\ 33')$&$-11.92$
\end{tabular}\endgroup
\end{center}
agree almost perfectly], he at once determined the following elements.
\begin{center}
\begingroup\small\setlength{\tabcolsep}{3pt}
\begin{tabular}{l r}
Epoch 1801 for Seeberg $\ .\ .$&$77^0\ 19'\ 38''4$\\
Daily tropical motion $\ .$&$770''764$\\
Period of revolution $\ .\ .\ .\ .\ .\ .\ .$&1681 days 11 hours\\
Eccentricity $\ .\ .\ .\ .\ .\ .\ .$&$0.0788132$\\[3pt]
Logarithm of the semimajor axis $\ .$&$0.4421085$\\
Aphelion 1801 $\ .\ .\ .\ .\ .\ .\ .\ .\ .$&$326^0\ 33'\ 10''$\\
Node 1801 $\ .\ .\ .\ .\ .\ .\ .\ .\ .\ .$&$80\ 54\ 52$\\
Inclination — $\ .\ .\ .\ .\ .\ .\ .\ .\ .\ .$&$10\ 37\ 48$
\end{tabular}\endgroup
\end{center}

With these elements Dr.\ \textsc{Gauss} then calculated the perturbations of $\CeresSymbol$ by $\JupiterSymbol$ anew, as follows.
\begin{center}\textit{Secular equations.}\end{center}
\begin{center}
\begingroup\small\setlength{\tabcolsep}{3pt}
\begin{tabular}{l r}
Annual advance of the aphelion relative to the equinoxes&$+121''3$\\
Annual advance of the nodes relative to the equinoxes $\ .\ .$&$+0.8$\\
Annual decrease in eccentricity $\ .\ .\ .\ .\ .\ .\ .\ .\ .$&$0.000005829$\\
Annual decrease in orbital inclination $\ .\ .\ .\ .\ .\ .\ .$&$+0''38$
\end{tabular}\endgroup
\end{center}
\Needspace{19\baselineskip}
\begin{center}\textit{Periodic equations.}\par\medskip
\textbf{In longitude.}\end{center}
\par\noindent\begin{minipage}[t]{.48\linewidth}\small
\[
\begin{aligned}
&-231''00\sin(\CeresSymbol-\JupiterSymbol)\\
&+492.37\sin2(\CeresSymbol-\JupiterSymbol)\\
&+43.85\sin3(\CeresSymbol-\JupiterSymbol)\\
&+10.00\sin4(\CeresSymbol-\JupiterSymbol)\\
&+3.03\sin5(\CeresSymbol-\JupiterSymbol)\\
&+1.04\sin6(\CeresSymbol-\JupiterSymbol)\\
&+0.40\sin7(\CeresSymbol-\JupiterSymbol)\\
&+0.16\sin8(\CeresSymbol-\JupiterSymbol)\\
&-58.90\sin(\JupiterSymbol+16^0\ 56'\ 37'')\\
&-598.69\sin(\CeresSymbol-2\JupiterSymbol-26^0\ 6'\ 58'')\\
&-437.75\sin(2\CeresSymbol-3\JupiterSymbol-10^0\ 54'\ 25'')
\end{aligned}
\]
\end{minipage}\hfill\begin{minipage}[t]{.515\linewidth}\small
\[
\begin{aligned}
&+54''90\sin(3\CeresSymbol-4\JupiterSymbol-10^0\ 32'\ 23'')\\
&+10.28\sin(4\CeresSymbol-5\JupiterSymbol-10\quad21\quad35'')\\
&+3.19\sin(5\CeresSymbol-6\JupiterSymbol-10\quad15\quad2'')\\
&+1.20\sin(6\CeresSymbol-7\JupiterSymbol-10\quad10\quad20'')\\
&+0.50\sin(7\CeresSymbol-8\JupiterSymbol-10\quad6\quad24'')\\
&+22.72\sin(2\CeresSymbol-\JupiterSymbol+36\quad0\quad12'')\\
&-51.80\sin(3\CeresSymbol-2\JupiterSymbol+32\quad47\quad57'')\\
&-5.74\sin(4\CeresSymbol-3\JupiterSymbol+30\quad41\quad32'')\\
&-1.65\sin(5\CeresSymbol-4\JupiterSymbol+28\quad53\quad58'')\\
&-0.69\sin(6\CeresSymbol-5\JupiterSymbol+28\quad13\quad52'')\\
&-0.26\sin(7\CeresSymbol-6\JupiterSymbol+26\quad41\quad16'')
\end{aligned}
\]
\end{minipage}\par
\SourceLeaf{234}%
\Needspace{19\baselineskip}
\begin{center}\textbf{In the radius vector.}\end{center}
\par\noindent\begin{minipage}[t]{.48\linewidth}\small
\[
\begin{aligned}
&-0.0000943\\
&+0.0010261\cos(\CeresSymbol-\JupiterSymbol)\\
&-0.0037679\cos2(\CeresSymbol-\JupiterSymbol)\\
&-0.0004176\cos3(\CeresSymbol-\JupiterSymbol)\\
&-0.0001069\cos4(\CeresSymbol-\JupiterSymbol)\\
&-0.0000347\cos5(\CeresSymbol-\JupiterSymbol)\\
&-0.0000127\cos6(\CeresSymbol-\JupiterSymbol)\\
&-0.0000050\cos7(\CeresSymbol-\JupiterSymbol)\\
&-0.0000021\cos8(\CeresSymbol-\JupiterSymbol)\\
&-0.0000614\cos(\CeresSymbol+24^0\ 13'\ 37'')\\
&+0.0002398\cos(\JupiterSymbol+22\quad58\quad33'')\\
&+0.0008325\cos(\CeresSymbol-2\JupiterSymbol-23^0\ 54'\ 51'')
\end{aligned}
\]
\end{minipage}\hfill\begin{minipage}[t]{.515\linewidth}\small
\[
\begin{aligned}
&+0.0025560\cos(2\CeresSymbol-3\JupiterSymbol-10^0\ 39'\ 20'')\\
&-0.0004698\cos(3\CeresSymbol-4\JupiterSymbol-10\quad51\quad24'')\\
&-0.0001048\cos(4\CeresSymbol-5\JupiterSymbol-10\quad46\quad51'')\\
&-0.0000358\cos(5\CeresSymbol-6\JupiterSymbol-10\quad40\quad14'')\\
&-0.0000142\cos(6\CeresSymbol-7\JupiterSymbol-10\quad34\quad4'')\\
&-0.0000061\cos(7\CeresSymbol-8\JupiterSymbol-10\quad28\quad56'')\\
&-0.0001313\cos(2\CeresSymbol-\JupiterSymbol+37\quad11\quad18'')\\
&+0.0002928\cos(3\CeresSymbol-2\JupiterSymbol+32\quad17\quad53'')\\
&+0.0000397\cos(4\CeresSymbol-3\JupiterSymbol+29\quad3\quad25'')\\
&+0.0000135\cos(5\CeresSymbol-4\JupiterSymbol+26\quad58\quad26'')\\
&+0.0000057\cos(6\CeresSymbol-5\JupiterSymbol+25\quad46\quad37'')\\
&+0.0000026\cos(7\CeresSymbol-6\JupiterSymbol+25\quad2\quad49'')
\end{aligned}
\]
\end{minipage}\par
\begin{center}\textbf{In latitude.}\end{center}
\par\noindent\begin{minipage}[t]{.48\linewidth}\small
\[
\begin{aligned}
&-11''60\sin(\JupiterSymbol-78^0\ 31'\ 50'')\\
&+14.29\sin(\CeresSymbol-2\JupiterSymbol+78^0\ 31'\ 50'')\\
&+28.73\sin(2\CeresSymbol-3\JupiterSymbol+78\quad31\quad50'')\\
&-4.81\sin(3\CeresSymbol-4\JupiterSymbol+78\quad31\quad50'')\\
&-0.96\sin(4\CeresSymbol-5\JupiterSymbol+78\quad31\quad50'')
\end{aligned}
\]
\end{minipage}\hfill\begin{minipage}[t]{.515\linewidth}\small
\[
\begin{aligned}
&-0''29\sin(5\CeresSymbol-6\JupiterSymbol+78^0\ 31'\ 50'')\\
&+5.43\sin(2\CeresSymbol-\JupiterSymbol-78\quad31\quad50'')\\
&+1.34\sin(3\CeresSymbol-2\JupiterSymbol-78\quad31\quad50'')\\
&+0.44\sin(4\CeresSymbol-3\JupiterSymbol-78\quad31\quad50'')
\end{aligned}
\]
\end{minipage}\par
Finally, Dr.\ \textsc{Gauss} used these formulas to recalculate the numerical values of the perturbations for the underlying observations. He thus found the following orbital elements, which we will henceforth denote by (VIII).
\begin{center}
\begin{tabular}{r|r|r|r}
\multicolumn{2}{c|}{Epoch for Seeberg}&Aphelion&Node\\\hline
1801&$77^0\ 19'\ 34''9$&$326^0\ 33'\ 37''$&$80^0\ 54'\ 59''$\\
1802&$155\ 28\ 35.1$&$326\ 35\ 39$&$80\ 55\ 0$\\
1803&$233\ 37\ 35.3$&$326\ 37\ 40$&$80\ 55\ 1$
\end{tabular}
\end{center}
\begin{center}
\begingroup\small\setlength{\tabcolsep}{3pt}
\begin{tabular}{l r}
Daily tropical motion $\ .\ .\ .\ .\ .\ .\ .\ .\ .\ .\ .$&$770''7951$\\
Tropical period of revolution $\ .\ .\ .\ .\ .\ .\ .\ .\ .\ .\ .\ .\ .$&1681 days 9 hours\\
Eccentricity (1801) $\ .\ .\ .\ .\ .\ .\ .\ .\ .\ .\ .\ .\ .\ .$&$0.0788352$\\
Inclination (1801) $\ .\ .\ .\ .\ .\ .\ .\ .\ .\ .\ .\ .\ .\ .\ .\ .$&$10^0\ 37'\ 56''0$\\
Logarithm of the semiaxis $\ .\ .\ .\ .\ .\ .\ .\ .\ .\ .$&$0.4420971$
\end{tabular}\endgroup
\end{center}
These new elements differ so little from the preceding ones that a fresh calculation of the perturbations from them can be regarded as superfluous, since the results cannot differ appreciably from those given above. Out of curiosity, Dr.\ \textsc{Gauss} tried to see how much the new elements VIII would differ from VII around the time of the encounter next year. Leaving aberration out of account, he finds the following for 28 June 1803 at 12 o’clock mean Seeberg time.
\begin{center}
\begin{tabular}{c|r|r}
\shortstack{According to\\elements}&Longitude of $\CeresSymbol$&\shortstack{Latitude of $\CeresSymbol$\\south}\\\hline
VII&$280^0\ 7'\ 43''$&$5^0\ 1'\ 21''$\\
VIII&$280\ 17\ 58$&$5\ 4\ 53$\\\hline
Difference&$10\ 15$&$3\ 32$
\end{tabular}
\end{center}
\SourceLeaf{235}%

Dr.\ \textsc{Gauss} closes his calculation with this remark. “Given the great trouble it takes to calculate a position of the planet even from the perturbation equations above, it seems hardly advisable as yet to add the equations for Saturn and those for Jupiter that depend on the squares and products of the eccentricities. Their inclusion would otherwise pose no difficulty beyond the \textit{Taedium} of mechanical calculation, especially since, after nearly half a year, we will now be able to give the elements a much higher degree of accuracy with the help of new observations. Since it is so laborious to calculate the value of the equations in each case, I think it will be important to consider how to shorten and simplify the tables for them. I will present my ideas on this to you in due course, along with a \textit{Specimen}.”
\begin{flushright}[\textsc{von Zach.}]\end{flushright}

